\chapter{Introduction}

\section{Introduction}

First responders operating in infrastructure-denied environments—burning structures, collapsed buildings, remote terrain, or contested urban zones—depend on communication and monitoring systems designed for normal conditions. When those conditions collapse, so do the systems. The 9/11 World Trade Center communications failure, in which 343 firefighters perished partly because their radios were incompatible with police channels \cite{NIST_9_11_comms}, and Hurricane Katrina in 2005, where the loss of cell towers paralysed coordinated rescue for 60,000 stranded survivors \cite{AFCEA_Katrina_comms}, illustrate a persistent architectural failure: centralised, infrastructure-dependent designs are most vulnerable precisely when they are most needed.

This project proposes a fundamentally different architecture. Each responder carries a self-contained wearable node built on the ESP32-C3 microcontroller and Semtech SX1276 LoRa transceiver that participates in a peer-to-peer mesh network requiring no pre-existing infrastructure. Every node runs a lightweight SHA-256 Proof-of-Authority (PoA) blockchain that creates a cryptographically verifiable, tamper-evident log of sensor events. A Federated Averaging (FedAvg) layer trains an on-device logistic regression anomaly detector collaboratively across all nodes by exchanging gradient weight vectors rather than raw biometric data, protecting physiological privacy while enabling collective learning.

The technical novelty lies in the co-design of all three layers—LoRa mesh, on-device blockchain, and over-the-air federated learning—on commodity microcontroller hardware under strict power, memory (320~KB SRAM), and bandwidth constraints (37.5~kbps maximum physical-layer throughput). The system is validated against a hybrid testbed combining physical ESP32-C3 nodes with software-simulated virtual nodes connected via an MQTT bridge, enabling performance characterisation at network scales impractical to instrument purely in hardware.

\section{Problem Statement}

The fundamental problem is the absence of a low-cost, infrastructure-independent, tamper-resistant wearable network that can provide real-time physiological monitoring, secure data logging, and distributed anomaly detection for first-responder teams where conventional communications infrastructure is unavailable.

Three distinct gaps drive this work:

\textbf{(1) Infrastructure dependency.} Existing wearable platforms conforming to NIST IR 8235 \cite{NIST_IR_8235} rely on cellular data or proprietary radio protocols requiring fixed base stations. When infrastructure is destroyed—exactly when monitoring is most critical—monitoring capability is lost.

\textbf{(2) Data integrity.} Sensor data stored on centralised servers is vulnerable to corruption and tampering. In forensic contexts (determining whether a responder's vitals were safe before entering a collapse zone), integrity must be provable without relying on any single custodian. Existing blockchain implementations target high-bandwidth, unconstrained environments; running one on a microcontroller with 320~KB SRAM over a 37.5~kbps radio requires a purpose-designed approach \cite{Semtech_SX1276_datasheet}.

\textbf{(3) Biometric privacy.} Heart rate, SpO2, and fall events are sensitive health data. Centralising them creates a single point of failure for privacy. Federated Learning ensures raw sensor readings never leave the node, but transmitting weight vectors over the low-bandwidth, variable-latency LoRa channel introduces engineering constraints that existing WiFi-centric FL frameworks do not address.

\section{Motivation}

The International Labour Organization estimates 2.3 million occupational fatalities annually, with emergency workers disproportionately represented \cite{ILO_safety_stats}. The NFPA reported 136 on-duty firefighter deaths in the United States in 2021 alone \cite{NFPA_2021_firefighter}; post-incident reviews consistently cite delayed physiological awareness as a contributing factor. Commercial PASS devices detect motion cessation and emit an audible alarm but provide no physiological data, no network communication, and no position fix \cite{IEEE_PST_wearables}.

Three concurrent technology developments now make the proposed architecture feasible. First, the SX1276 LoRa transceiver achieves link budgets exceeding 155~dB with sub-100~mW power, enabling reliable links at 200--1000~m in urban and forested environments \cite{Semtech_SX1276_datasheet}. Second, peer-to-peer LoRa mesh has been demonstrated for mountain rescue and industrial safety without any gateway dependency \cite{LoRa_mesh_rescue_2024}. Third, the ESP32-C3 executes SHA-256 via mbedTLS in under 2~ms per block, making on-device blockchain viable within a 5-second sensor reporting cycle.

Existing alternatives are inadequate. FirstNet, the US public safety LTE network, loses coverage when cell towers are destroyed \cite{Preprints_FirstNet_5G_review}—as demonstrated in the 2019--2020 Australian bushfires where large areas lost all cellular connectivity for days. Iridium and Globalstar satellite options carry prohibitive per-device cost and latency. P25 and TETRA land mobile radio systems do not carry physiological sensor data or provide cryptographic data integrity. The FL motivation is similarly specific: centralised aggregation of health data is vulnerable to gradient inversion attacks even when data is encrypted in transit \cite{FL_privacy_threats_survey_2024}, while FedAvg with differential privacy provides a rigorous bound on per-individual information leakage \cite{FL_DP_McMahan2018}.

\section{Objectives}

\begin{enumerate}
    \item Integrate MPU6050 IMU, MAX30102 PPG sensor, and SX1276 LoRa transceiver on the ESP32-C3 for concurrent fall detection, heart rate, and SpO2 measurement.
    \item Implement SHA-256 PoA blockchain directly on the ESP32-C3 within the 320~KB SRAM constraint.
    \item Establish reliable bidirectional LoRa P2P communication between nodes with per-block hash verification.
    \item Implement FedAvg with Gaussian DP noise over LoRa and demonstrate weight convergence in a bounded number of rounds.
    \item Implement two-phase fall detection and SpO2/HR anomaly alerting with blockchain-logged, LoRa-propagated alert blocks.
    \item Develop a Python simulation (Gudmundson channel model, blockchain, FL) for up to 20 nodes and validate against hardware measurements via MQTT bridge.
\end{enumerate}

\section{Brief Methodology of the Project}

The project follows a milestone-driven methodology: five software simulation milestones (M1--M5) completed prior to hardware procurement, followed by six hardware milestones (M6--M11) building toward full system integration.

\textbf{Simulation phase (M1--M5):} M1 models the SX1276 physical layer with Gudmundson log-normal shadowing to replicate empirical PDR curves. M2 implements a Python SHA-256 hash chain providing ground-truth hashes for firmware verification. M3 processes UCI HAR accelerometer data \cite{UCI_HAR_dataset} for fall detection algorithm development. M4 implements FedAvg with Gaussian DP noise, demonstrating L2-norm convergence across simulated nodes. M5 exercises the full system at 5, 10, and 20 nodes, characterising blockchain propagation latency and FL convergence rate as a function of network size.

\textbf{Hardware phase (M6--M11):} M6 verifies sensor initialisation and valid readings. M7 establishes bidirectional LoRa block exchange with SHA-256 hash verification. M8 exercises multi-block hash chain and prevHash linkage against the Python simulation. M9 implements FedAvg weight exchange on hardware with convergence demonstration. M10 tests the full operational loop including a 10-minute soak test for heap stability. M11 connects physical nodes to the Python simulation mesh via USB-serial JSON bridge and Mosquitto MQTT, verifying three-way chain hash agreement.

\section{Key Contributions}

\begin{itemize}
\item First simultaneous co-integration of SHA-256 PoA blockchain and FedAvg over a LoRa P2P radio on a commodity ESP32-C3 (320~KB SRAM).
\item Airtime-aware TDMA-scheduled FL protocol accounting for the 1.8-second channel occupancy of a 124-byte LoRa block at SF10/BW125.
\item Hybrid MQTT testbed enabling physical ESP32 nodes and Python-simulated virtual nodes to participate in the same mesh as first-class peers.
\end{itemize}

The remainder of this report is organised as follows: Chapter~2 surveys the literature and theoretical foundations; Chapter~3 presents system design; Chapter~4 covers implementation; Chapter~5 reports results; Chapter~6 concludes.
